This project uses a restricted subset of ISO C11 intended to improve memory safety, type safety, portability, analyzability, and maintainability.

The C11 standard is organized into the following areas:

\begin{itemize}
    \item general principles;
    \item language subset;
    \item formatting;
    \item naming;
    \item Variables;
    \item module boundaries;
    \item public and private visibility;
    \item type safety;
    \item initialization;
    \item control flow;
    \item functions;
    \item memory safety;
    \item arrays and strings;
    \item pointers;
    \item integer safety;
    \item error handling;
    \item preprocessor usage;
    \item external libraries;
    \item concurrency;
    \item documentation and verification.
\end{itemize}

\section{General Principles}

Clarity First: Code should be as readable as prose, spoken aloud.

Simplicity is Beauty: Short, understandable code reduces bugs; remove dead code and refactor duplicates into functions.

Consistent Style: Align with existing project style; use format‑conversion tools if needed.

\section{C11 Language Subset}
\subsubsection{Permitted Language}

Source code shall conform to ISO C11. Compiler extensions shall not be used unless they are isolated behind a documented platform abstraction.

Rules prohibit unrestricted control flow and recursion. Instead it requires statically bounded loops, restricted runtime allocation, requires return-value checks, limit the preprocessor use, restrict the use of pointers, and require warning-free compilation and static analysis.
\subsubsection{Prohibited or Restricted Features}

The following features require documented justification:

\begin{itemize}
	\item unions;
	\item bit-fields;
	\item flexible array members;
	\item volatile objects;
	\item atomics;
	\item function pointers;
	\item implementation-defined behavior;
	\item compiler extensions;
	\item direct hardware access;
	\item dynamic memory allocation.
	
\end{itemize}

The following features are prohibited:

\begin{itemize}
	\item \texttt{goto};
	\item \texttt{setjmp};
	\item \texttt{longjmp};
	\item direct recursion;
	\item indirect recursion;
	\item variable-length arrays;
	\item implicit function declarations;
	\item implicit integer-to-pointer conversions;
	\item implicit pointer-to-integer conversions;
	\item function-like macros where an inline function is practical;
	\item compiler-specific types in portable modules;
	\item unbounded loops in deterministic modules;
	\item unrestricted global mutable state.
\end{itemize}


\section{Formatting}
\subsubsection{Indentation}

Tabs shall be used for indentation. Spaces shall not be used to represent indentation.

One tab shall represent one logical indentation level. The display width of a tab shall be four spaces in the project editor configuration.

\subsubsection{Alignment}

Spaces may be used for alignment within a line, but alignment shall not create a requirement for additional whitespace on later lines.

\subsubsection{Braces}

Braces shall always be used for compound statements, including statements containing only one instruction.

Opening braces shall appear on the same line as the associated statement.

\begin{lstlisting}[language=C]
if (condition) {
	process_item();
} else {
	handle_error();
}
\end{lstlisting}

\subsubsection{Whitespace}

There shall be spaces:

\begin{itemize}
	\item around binary operators;
	\item after commas;
	\item between a type and a variable name where appropriate.
\end{itemize}

After every keyword (if, while, etc) of the C language, a space must be added. e.g.
\begin{lstlisting}[language=C]
if (i == 1)
while (i == 0)
\end{lstlisting}
keywords include:
\begin{itemize}
	\item if
	\item switch
	\item case
	\item for
	\item do
	\item while
\end{itemize}
but not:
\begin{itemize}
	\item sizeof
	\item typeof
	\item alignof
	\item attribute
\end{itemize}

There shall be no trailing whitespace.

\subsubsection{Line Length}

Lines should not exceed 100 characters. A longer line is permitted only when breaking it would reduce readability.

\subsubsection{Automated Formatting}

All C files shall be processed using the project-approved formatter configuration. Manual formatting preferences shall not override the project formatter.

Package indent under Debian and Ubuntu Linux distributions is recommended for formatting C code.
On windows it can be found here: https://gnuwin32.sourceforge.net/packages/indent.htm

\section{Naming}
\subsection{General Naming Requirements}

Names shall describe the purpose, behavior, or ownership of the object they identify. Abbreviations shall be avoided unless they are widely understood within the project.

Names shall:

\begin{itemize}
	\item use lowercase \texttt{snake\_case} for functions and variables;
	\item use uppercase \texttt{UPPER\_SNAKE\_CASE} for macros;
	\item use descriptive names for types;
	\item distinguish counts, sizes, indexes, capacities, and offsets;
	\item avoid names that differ only by capitalization;
	\item avoid names that shadow variables in an enclosing scope;
	\item avoid single-letter names except for short loop indexes;
	\item avoid misleading names.
\end{itemize}

\subsection{Functions}

Functions shall use verb-based names that describe their operation.

\begin{lstlisting}[language=C]
buffer_init();
buffer_write();
buffer_clear();
packet_is_valid();
\end{lstlisting}

Boolean-returning functions should use names such as:

\begin{lstlisting}[language=C]
is_valid();
has_error();
can_write();
\end{lstlisting}

\subsection{Variables}

Variables shall describe their content and, where useful, their unit.

\begin{lstlisting}[language=C]
size_t packet_length;
uint32_t timeout_ms;
uint8_t retry_count;
bool connection_is_open;
\end{lstlisting}

Names such as \texttt{data}, \texttt{value}, \texttt{temp}, and \texttt{result} should be avoided unless their meaning is clear from the local context.

\subsection{Types}

Project-defined types shall use a consistent project convention.

\begin{lstlisting}[language=C]
typedef struct buffer buffer_t;
typedef enum device_state device_state_t;
\end{lstlisting}

Structure tags and typedef names shall not unnecessarily duplicate one another.

\subsection{Constants and Macros}

Constants should preferably be represented using typed constants, enumerations, or \texttt{static const} objects instead of macros.

\begin{lstlisting}[language=C]
static const size_t max_packet_length = 1024U;
\end{lstlisting}

Macros shall use uppercase names with a project-specific prefix.

\begin{lstlisting}[language=C]
#define PROJECT_MAX_PACKET_LENGTH 1024U
\end{lstlisting}

\subsection{Private Names}

Private file-scope functions and variables shall use \texttt{static}. A project prefix may be used for private names when required to avoid collisions with external symbols.


\section{Varibles}
\subsubsection{General Varible Requirements}

Group local variables together by type.

Do not declare variables after first executable statement.

Each variable has a single purpose; avoid globals.

Prevent name clashes, respect byte order in communication structs, and never use uninitialized variables.

\section{Module Boundaries}
\subsubsection{Module Structure}

Each module shall consist of:

\begin{itemize}
	\item one public header file, when an external interface is required;
	\item one source file;
	\item private declarations contained in the source file;
	\item unit tests for externally visible behavior.
\end{itemize}

A module shall have one primary responsibility.

\subsubsection{Header Independence}

Every public header shall compile independently when included into an otherwise empty translation unit.

\begin{lstlisting}[language=C]
#include "buffer.h"

int main(void)
{
	return 0;
}
\end{lstlisting}

\subsubsection{Include Ownership}

A source file shall include the header that declares its public interface before other project headers.

\begin{lstlisting}[language=C]
#include "buffer.h"

#include <stddef.h>
#include <stdint.h>

#include "memory.h"
\end{lstlisting}

\subsubsection{Circular Dependencies}

Circular dependencies between modules are prohibited.

Modules shall depend on abstractions or narrow interfaces rather than including unrelated internal headers.


\section{Visibility}
\subsection{Public Declarations}

A public declaration is any declaration exposed through a module's interface header.

Public declarations shall be limited to:

\begin{itemize}
	\item functions intended for use by other modules;
	\item opaque object types;
	\item public enumerations;
	\item public constants;
	\item public error codes;
	\item public configuration structures when required by the interface.
\end{itemize}

\subsection{Private Declarations}

Private functions, variables, constants, and types shall be defined in the corresponding source file.

Private file-scope declarations shall use the \texttt{static} specifier.

\begin{lstlisting}[language=C]
static uint32_t calculate_checksum(const uint8_t *data,
                                   size_t length);

static uint32_t checksum_seed = 0U;
\end{lstlisting}

\subsection{Global Variables}

Global mutable variables are prohibited.

File-scope variables shall be private unless there is a documented architectural reason for making them public.

Public variables shall not be used for ordinary module communication. Functions shall be used instead.

\subsection{Opaque Types}

Public headers should expose opaque structures rather than structure definitions.

\begin{lstlisting}[language=C]
/* buffer.h */

#ifndef PROJECT_BUFFER_H
#define PROJECT_BUFFER_H

#include <stddef.h>
#include <stdint.h>

typedef struct buffer buffer_t;

buffer_t *buffer_create(uint8_t *storage, size_t capacity);
int buffer_write(buffer_t *buffer,
                 const uint8_t *data,
                 size_t length);
size_t buffer_size(const buffer_t *buffer);

#endif
\end{lstlisting}

The structure definition remains private:

\begin{lstlisting}[language=C]
/* buffer.c */

struct buffer {
	uint8_t *storage;
	size_t capacity;
	size_t length;
};
\end{lstlisting}


\section{Type Safety}
\subsection{Fixed-Width Integer Types}

Fixed-width integer types shall be used when the width or signedness is important.

\begin{lstlisting}[language=C]
#include <stdint.h>

uint32_t packet_length;
int32_t temperature;
uint8_t status;
\end{lstlisting}

The types \texttt{int8\_t}, \texttt{uint8\_t}, and similar types shall be used instead of assuming the size of \texttt{char}, \texttt{short}, or \texttt{int}.

\subsection{Boolean Values}

Boolean values shall use \texttt{bool} from \texttt{<stdbool.h>}.

\begin{lstlisting}[language=C]
#include <stdbool.h>

bool packet_is_valid;
\end{lstlisting}

Integer values shall not be used as booleans unless the conversion is explicit and documented.

\subsection{Enumerations}

Enumerations shall be used for finite sets of named states.

\begin{lstlisting}[language=C]
typedef enum {
	MODE_IDLE = 0,
	MODE_ACTIVE,
	MODE_ERROR
} mode_t;
\end{lstlisting}

\subsection{Explicit Conversions}

Conversions that may change signedness, width, precision, or interpretation shall be explicit.

\begin{lstlisting}[language=C]
uint32_t value = 0U;
uint16_t reduced_value = (uint16_t)value;
\end{lstlisting}

The conversion shall be safe. A cast shall not be used merely to silence a compiler warning.

\subsection{Const Correctness}

The \texttt{const} qualifier shall be used whenever an object is not modified through a pointer.

\begin{lstlisting}[language=C]
void process_packet(const uint8_t *data, size_t length);
\end{lstlisting}

\subsection{Additional Rules}

\begin{itemize}
	\item Do not compare signed and unsigned values without an explicit conversion or documented reason.
	\item Do not use plain \texttt{char} for binary data.
	\item Do not use \texttt{float} or \texttt{double} for exact quantities.
	\item Do not rely on the size of \texttt{int}.
	\item Do not use casts to convert unrelated pointer types.
	\item Do not use type-punning through incompatible pointer types.
	\item Do not use an enumeration outside its defined semantic domain.
\end{itemize}


\section{Initialization}
\subsubsection{General Requirement}

Every object shall be initialized before it is read.

Initialization shall occur as close as practical to the declaration.

\begin{lstlisting}[language=C]
uint32_t retry_count = 0U;
bool packet_is_valid = false;
\end{lstlisting}

Uninitialized local variables are prohibited.

\subsubsection{Default Initialization}

Objects shall be initialized to a valid and documented default state.

\begin{lstlisting}[language=C]
struct device device = {
	.state = DEVICE_STATE_OFF,
	.timeout_ms = 1000U,
	.error_code = ERROR_NONE
};
\end{lstlisting}

Designated initializers shall be used for structures when they improve clarity or protect against changes in member ordering.

\subsubsection{Structures}

Structure members shall be initialized explicitly when their zero value is not the required default.

The following form is preferred for complete initialization:

\begin{lstlisting}[language=C]
struct configuration configuration = {
	.enabled = true,
	.retry_limit = 3U,
	.timeout_ms = 500U
};
\end{lstlisting}

\subsubsection{Arrays}

Arrays shall be initialized explicitly when partial initialization could hide an error.

\begin{lstlisting}[language=C]
uint8_t header[HEADER_LENGTH] = {0U};
\end{lstlisting}

\subsubsection{Pointers}

Pointers shall be initialized to a valid object or to \texttt{NULL}.

\begin{lstlisting}[language=C]
uint8_t *buffer = NULL;
\end{lstlisting}

A pointer shall not be dereferenced unless its validity has been established.

\subsubsection{Static Storage}

Static and file-scope objects shall have an explicit initializer when the initial state is important to the module contract.

\subsubsection{Initialization Functions}

Initialization functions shall:

\begin{itemize}
	\item be called before any other operation on the object;
	\item be safe to call only according to their documented contract;
	\item establish all required invariants;
	\item report failure explicitly;
	\item leave the object in a documented state when initialization fails.
\end{itemize}

\begin{lstlisting}[language=C]
int device_init(device_t *device);
\end{lstlisting}

\subsubsection{Partial Initialization}

An object shall not be exposed to other modules until initialization is complete.

If initialization requires multiple steps, cleanup shall be defined for each partially initialized state.


\section{Control Flow}
\subsection{Control Flow}

The following are prohibited:

\begin{itemize}
	\item \texttt{goto};
	\item \texttt{setjmp};
	\item \texttt{longjmp};
	\item recursion;
	\item deeply nested conditionals;
	\item hidden control flow in macros.
\end{itemize}

\subsection{Loops}

Loops shall have a clear termination condition.

Loops in deterministic or safety-critical code shall have a statically reviewable upper bound.

\begin{lstlisting}[language=C]
for (size_t index = 0U; index < MAX_ITEMS; ++index) {
	process_item(items[index]);
}
\end{lstlisting}

A loop controlled by external input shall still have a maximum iteration limit.


\section{Functions}
\subsection{Function Requirements}

Functions shall:

\begin{itemize}
	\item have one primary responsibility;
	\item have a clearly defined contract;
	\item validate input parameters;
	\item return an explicit status for recoverable failure;
	\item avoid hidden side effects;
	\item avoid modifying caller data unless documented;
	\item remain short enough for complete review.
	\item Limit nesting depth to four levels.
\end{itemize}

Functions should normally have no more than 10 local variables. If more are required, the function should be refactored into smaller functions.

Every function that may be accessed from outside its module MUST include a function prototype (declaration).

Function declarations and definitions that accept no parameters MUST explicitly use (void), never empty parentheses ().

Function implementations MUST place the return type and any storage-class or qualifier keywords on a separate line from the function name.

A function SHOULD have at most 2 return points: one early return at the top of the function for parameter/argument validation, and one return at the end of the function. Avoid return statements in the middle of the function body; use a return-value variable to accumulate the result instead.



\subsection{Parameter Validation}

Public functions shall validate all pointer parameters, lengths, ranges, and state assumptions that are not guaranteed by the type system.

\subsection{Return Values}

The caller shall check every non-void return value.

If a return value is intentionally ignored, the reason shall be explicit.

\begin{lstlisting}[language=C]
(void)log_message("startup complete");
\end{lstlisting}


\section{Memory Safety}
\subsubsection{General Requirement}

Code shall not produce:

\begin{itemize}
	\item out-of-bounds accesses;
	\item use-after-free;
	\item double-free;
	\item invalid frees;
	\item null-pointer dereferences;
	\item uninitialized reads;
	\item memory leaks;
	\item access through a dangling pointer;
	\item writes through a read-only object;
	\item integer-overflow-derived allocation sizes.
\end{itemize}

For maximum saftey general rules include:

\begin{itemize}
    \item no heap allocation in core logic;
    \item no heap allocation unless needed;
    \item caller-provided buffers where possible;
    \item explicit buffer lengths;
    \item no hidden ownership transfer;
    \item no ownership implied only by naming;
    \item cleanup paths tested.
\end{itemize}

\subsubsection{Allocation Policy}

Static allocation and caller-provided storage are preferred.

Dynamic allocation shall be avoided in ordinary modules and is prohibited after initialization in deterministic or safety-critical modules.

Every allocation shall have:

\begin{itemize}
	\item one owner;
	\item one documented release path;
	\item a checked failure path;
	\item a defined maximum size;
	\item a test for boundary conditions.
\end{itemize}

\subsubsection{Allocation Arithmetic}

Allocation-size arithmetic shall be checked before allocation.

\begin{lstlisting}[language=C]
if (count > SIZE_MAX / sizeof(*items)) {
	return ERROR_SIZE_OVERFLOW;
}

items = malloc(count * sizeof(*items));
if (items == NULL) {
	return ERROR_OUT_OF_MEMORY;
}
\end{lstlisting}

\subsubsection{Ownership}

Functions shall document whether pointer parameters are:

\begin{itemize}
	\item borrowed and not retained;
	\item owned by the caller;
	\item transferred to the callee;
	\item returned with ownership transferred to the caller.
\end{itemize}


\section{Arrays and Strings}
\subsubsection{Array Bounds}

Every array access shall be demonstrably within bounds.

A pointer to an array shall be accompanied by its length unless the length is guaranteed by the interface contract.

\begin{lstlisting}[language=C]
int process_bytes(const uint8_t *data, size_t length);
\end{lstlisting}

\subsubsection{String Handling}

The following functions are prohibited:

\begin{itemize}
	\item \texttt{gets};
	\item \texttt{strcpy};
	\item \texttt{strcat};
	\item \texttt{sprintf};
	\item \texttt{scanf} without a field-width limit.
\end{itemize}

String operations shall specify destination capacity and shall preserve null termination.

\subsubsection{Binary Data}

Binary buffers shall use \texttt{uint8\_t} and an explicit length. They shall not be treated as null-terminated strings.


\section{Pointers}
\input{c/cPointers.tex}

\section{Integer Safety}
\input{c/cIntSafety.tex}

\section{Error Handling}
\input{c/cErrorHandling.tex}

\section{Preprocessor}
\input{c/cPreprocessor.tex}

\section{External Libraries}
\subsubsection{Library Policy}

The project shall depend only on:

\begin{itemize}
	\item the ISO C11 standard library;
	\item operating-system interfaces required by the target;
	\item project-owned libraries;
	\item third-party libraries approved in the dependency register.
\end{itemize}

A third-party library shall not be added solely for convenience when the required functionality can be implemented safely using C11 or existing project code.

\subsubsection{Approval Requirements}

Each external library shall have:

\begin{itemize}
	\item a documented purpose;
	\item a version pinned in the build system;
	\item a license record;
	\item a maintenance status;
	\item a security review;
	\item a documented public API surface;
	\item a removal or replacement plan where practical.
\end{itemize}

\subsubsection{Wrapper Requirement}

External libraries shall be accessed through a project-owned wrapper module.

Application modules shall not include third-party headers directly unless the dependency is explicitly approved as part of the public architecture.

\begin{lstlisting}[language=C]
/* project_json.h */

#ifndef PROJECT_JSON_H
#define PROJECT_JSON_H

int project_json_parse(const char *text);

#endif
\end{lstlisting}


\section{Concurrency}
\subsection{Applicability}

Concurrency includes operating-system threads, interrupt handlers, callbacks, signal handlers, asynchronous tasks, and hardware event handlers.

Concurrency shall not be introduced unless it is required by the system design.

\subsection{Shared State}

Shared mutable state shall be minimized.

Every shared object shall have a documented synchronization strategy.

The following are prohibited:

\begin{itemize}
	\item undocumented shared mutable state;
	\item data races;
	\item unsynchronized access to shared objects;
	\item relying on timing for correctness;
	\item using \texttt{volatile} as a replacement for synchronization.
\end{itemize}

\subsection{Ownership}

Each shared object shall have one of the following documented policies:

\begin{itemize}
	\item protected by a mutex;
	\item accessed only by one thread;
	\item accessed atomically;
	\item immutable after initialization;
	\item exchanged through a documented message-passing interface.
\end{itemize}

\subsection{Mutexes}

Mutex acquisition and release shall be visibly paired.

\begin{lstlisting}[language=C]
mutex_lock(&state_mutex);

update_shared_state();

mutex_unlock(&state_mutex);
\end{lstlisting}

A mutex shall not be held while calling an external function unless the locking contract explicitly permits it.

\subsection{Lock Ordering}

When multiple locks are required, the project shall define a global lock ordering.

Locks shall always be acquired in that order to prevent deadlock.

\subsection{Atomics}

C11 atomic objects shall be used only when their memory-ordering behavior is understood and documented.

The default sequentially consistent ordering shall be preferred unless a different ordering is justified.

\subsection{Interrupts and Signal Handlers}

Signal handlers and interrupt handlers shall perform minimal work.

They shall not call functions that are not safe in their execution context.

Communication with the main application shall use an approved mechanism, such as an atomic flag, queue, or event object.

\subsection{Thread Cancellation}

Thread cancellation shall not be used unless cleanup behavior is fully defined.

Resources shall not be leaked when a thread exits or is cancelled.

\subsection{Testing}

Concurrent code shall be tested under:

\begin{itemize}
	\item repeated execution;
	\item increased scheduling variability;
	\item thread and lock analysis;
	\item race detection where supported;
	\item sanitizer-enabled builds where supported.
\end{itemize}


\section{Documentation}
\subsection{Documentation Requirements}

Documentation shall explain intent, contracts, assumptions, ownership, and limitations. It shall not merely repeat the implementation.

Documentation shall be updated whenever behavior or interfaces change.

\subsection{Public Interfaces}

Every public function shall document:

\begin{itemize}
	\item its purpose;
	\item each parameter;
	\item the return value;
	\item possible error conditions;
	\item ownership and lifetime rules;
	\item thread-safety requirements;
	\item whether pointers may be \texttt{NULL}.
\end{itemize}

\begin{lstlisting}[language=C]
/**
 * @brief Writes bytes to a buffer.
 *
 * @param buffer Buffer that receives the data.
 * @param data Input data. The memory is borrowed and is not retained.
 * @param length Number of bytes to write.
 *
 * @return 0 on success.
 * @return -1 if an argument is invalid.
 * @return -2 if insufficient capacity is available.
 *
 * @threadsafe No. The caller shall provide synchronization.
 */
int buffer_write(buffer_t *buffer,
                 const uint8_t *data,
                 size_t length);
\end{lstlisting}

\subsection{Module Documentation}

Each module shall document:

\begin{itemize}
	\item its responsibility;
	\item its public interface;
	\item its dependencies;
	\item its ownership model;
	\item its initialization requirements;
	\item its thread-safety model;
	\item its error-reporting strategy.
\end{itemize}

\subsection{Comments}

Comments shall explain why the code exists, not simply describe what the code visibly does.

Poor comment:

\begin{lstlisting}[language=C]
// Increment index
index++;
\end{lstlisting}

Useful comment:

\begin{lstlisting}[language=C]
// The hardware requires a delay after reset before the status register becomes valid.
delay_ms(HARDWARE_RESET_DELAY_MS);
\end{lstlisting}

For multi-line comments use space+asterisk for every line

\subsection{TODO Comments}

TODO comments shall include an issue identifier or owner.

\begin{lstlisting}[language=C]
// TODO(PROJ-142): Replace polling with event notification.
\end{lstlisting}

Unowned TODO comments shall not be committed to production code.

\subsection{Generated Documentation}

Public interfaces should use a consistent documentation format, such as Doxygen. Generated documentation shall be treated as a build artifact and shall not replace source-level documentation.


\section{Tooling}
\subsection{Compiler}

The compiler shall be configured for C11:

\begin{lstlisting}[language=bash]
-std=c11
\end{lstlisting}

The following warning categories should be enabled for GCC or Clang:

\begin{lstlisting}[language=bash]
-Wall
-Wextra
-Wpedantic
-Wconversion
-Wsign-conversion
-Wshadow
-Wcast-qual
-Wcast-align
-Wstrict-prototypes
-Wmissing-prototypes
-Wformat=2
-Wundef
-Werror
\end{lstlisting}

The exact warning set may vary by compiler, but all warnings shall be resolved or formally justified.

\subsection{Analysis}

The project shall use:

\begin{itemize}
	\item compiler warnings;
	\item a static analyzer;
	\item AddressSanitizer during testing;
	\item UndefinedBehaviorSanitizer during testing;
	\item a formatter;
	\item unit tests;
	\item boundary and negative tests.
\end{itemize}

\subsection{Required Build Checks}

A merge shall be rejected if:

\begin{itemize}
	\item formatting fails;
	\item compilation produces warnings;
	\item static analysis produces unreviewed findings;
	\item sanitizers detect an error;
	\item unit tests fail;
	\item coverage or safety requirements are not met.
\end{itemize}


\section{Exceptions}
\subsubsection{Purpose}

Exceptions provide a controlled method for deviating from this standard. They shall not be used to avoid normal engineering review.

\subsubsection{When an Exception Is Required}

An exception is required when code:

\begin{itemize}
	\item uses a prohibited language feature;
	\item introduces an unapproved external dependency;
	\item exposes a normally private symbol;
	\item uses dynamic allocation in a restricted module;
	\item suppresses a compiler or static-analysis warning;
	\item violates a mandatory formatting or naming rule;
	\item requires implementation-defined or compiler-specific behavior;
	\item cannot satisfy a memory, type, or concurrency rule.
\end{itemize}

\subsubsection{Exception Record}

Each exception shall document:

\begin{itemize}
	\item the rule identifier;
	\item the affected file, module, or function;
	\item the reason for the deviation;
	\item the risks introduced;
	\item the mitigation or compensating control;
	\item the approving reviewer;
	\item the date of approval;
	\item the planned review or expiration date.
\end{itemize}

\begin{longtable}{p{0.22\textwidth} p{0.68\textwidth}}
\toprule
Field & Value \\
\midrule
Rule identifier & C-MEM-004 \\
Affected module & \texttt{network\_buffer.c} \\
Deviation & Dynamic allocation is used after initialization. \\
Reason & Incoming packet size is not known during startup. \\
Mitigation & Maximum allocation is bounded and failure is handled. \\
Approver & Technical lead \\
Review date & 2026-10-03 \\
Status & Approved \\
\bottomrule
\end{longtable}

\subsubsection{Warning Suppression}

Warning suppression shall be localized to the smallest possible scope.

A suppression shall include a comment explaining:

\begin{itemize}
	\item which warning is suppressed;
	\item why it is safe;
	\item why the code cannot be rewritten to avoid it.
\end{itemize}

Suppressing warnings globally is prohibited.

\subsubsection{Review}

Exceptions shall be reviewed when the affected module changes and at each major release.

Expired or no-longer-needed exceptions shall be removed.
